\ifdefined\gkver\else\def\gkver{student}\fi
\documentclass[\gkver]{gaokaozhenti}
\newcommand{\ccnum}[1]{{\fontspec{Noto Serif CJK SC}#1}}
\begin{document}

\chapter{2010年全国理综Ⅱ}
%% paperType: 理综物理部分

\begin{enumerate}
\item
%% number: 14
%% paperName: 2010年全国理综Ⅱ第 14 题 6 分
%% typeId: 1
%% type: 单选题
%% chapter: 原子和原子核
%% point: 人工核反应
%% method: 无
%% score: 6
%% degree: 650
%% duplicateId: 0
%% body:
原子核 \ce{^{A}_{Z}X} 与氘核 \ce{^{2}_{1}H} 反应生成一个 $\alpha$ 粒子和一个质子。由此可知\xzanswer{D}
\fourchoices[answer=D]
{$A=2$，$Z=1$}
{$A=2$，$Z=2$}
{$A=3$，$Z=3$}
{$A=3$，$Z=2$}
%% answer: D
%% memo:
\memoanswer{
核反应方程为 $\ce{^{A}_{Z}X + ^{2}_{1}H -> ^{4}_{2}He + ^{1}_{1}H}$，应用质量数与电荷数的守恒 $A+2=4+1$，$Z+1=2+1$，解得 $A=3$，$Z=2$，选项 D 正确。
}

\item
%% number: 15
%% paperName: 2010年全国理综Ⅱ第 15 题 6 分
%% typeId: 2
%% type: 多选题
%% chapter: 机械波
%% point: 波的图象
%% method: 无
%% score: 6
%% degree: 650
%% duplicateId: 0
%% body:
一简谐横波以 $4\Ums$ 的波速沿 $x$ 轴正方向传播。已知 $t=0$ 时的波形如图所示，则\xzanswer{AB}
\onepicture[w=5.5cm]{figs/15.png}
\fourchoices[answer=AB]
{波的周期为 $1\Us$}
{$x=0$ 处的质点在 $t=0$ 时向 $y$ 轴负向运动}
{$x=0$ 处的质点在 $t=\frac{1}{4}\Us$ 时速度为 0}
{$x=0$ 处的质点在 $t=\frac{1}{4}\Us$ 时速度值最大}
%% answer: AB
%% memo:
\memoanswer{
由图可得半波长为 $2\Um$，波长为 $4\Um$。周期 $T=\frac{\lambda}{v}=\frac{4}{4}\Us=1\Us$，选项 A 正确。波沿 $x$ 轴正方向传播，则 $x=0$ 的质点在沿 $y$ 轴的负方向运动，选项 B 正确。$x=0$ 的质点位移为振幅的一半，要运动到平衡位置的时间是 $\frac{1}{3}\times\frac{T}{4}=\frac{1}{12}\Us$，则 $t=\frac{1}{4}\Us$ 时，$x=0$ 的质点不在平衡位置也不在最大位移处，故 CD 错误。
}

\item
%% number: 16
%% paperName: 2010年全国理综Ⅱ第 16 题 6 分
%% typeId: 2
%% type: 多选题
%% chapter: 气体性质
%% point: 热力学第一定律
%% method: 无
%% score: 6
%% degree: 450
%% duplicateId: 0
%% body:
如图，一绝热容器被隔板 K 隔开 a、b 两部分。已知 a 内有一定量的稀薄气体，b 内为真空，抽开隔板 K 后，a 内气体进入 b，最终达到平衡状态。在此过程中\xzanswer{BD}
\onepicture[w=3.5cm]{figs/16.png}
\fourchoices[answer=BD]
{气体对外界做功，内能减少}
{气体不做功，内能不变}
{气体压强变小，温度降低}
{气体压强变小，温度不变}
%% answer: BD
%% memo:
\memoanswer{
绝热容器内的稀薄气体与外界没有热交换，$Q=0$。稀薄气体向真空中扩散没有做功，$W=0$。根据热力学第一定律稀薄气体的内能不变，则温度不变。稀薄气体扩散体积增大，压强必减小。选项 BD 正确。
}

\item
%% number: 17
%% paperName: 2010年全国理综Ⅱ第 17 题 6 分
%% typeId: 1
%% type: 单选题
%% chapter: 电场
%% point: 带电粒子的平衡
%% method: 无
%% score: 6
%% degree: 650
%% duplicateId: 0
%% body:
在雷雨云下沿竖直方向的电场强度为 $10^{4}\UVm$。已知一半径为 $1\Umm$ 的雨滴在此电场中不会下落，取重力加速度大小为 $10\Umsq$，水的密度为 $10^{3}\Ukgmc$。这雨滴携带的电荷量的最小值约为\xzanswer{B}
\fourchoices[answer=B]
{$2\times10^{-9}\UC$}
{$4\times10^{-9}\UC$}
{$6\times10^{-9}\UC$}
{$8\times10^{-9}\UC$}
%% answer: B
%% memo:
\memoanswer{
带电雨滴在电场力和重力作用下保持静止，则 $mg=qE$，得：$q=\frac{mg}{E}=\frac{\rho Vg}{E}=\frac{\rho\times\frac{4}{3}\pi r^{3}g}{E}=4\times10^{-9}\UC$，选项 B 正确。
}

\item
%% number: 18
%% paperName: 2010年全国理综Ⅱ第 18 题 6 分
%% typeId: 1
%% type: 单选题
%% chapter: 电磁感应
%% point: 电磁感应受力分析
%% method: 无
%% score: 6
%% degree: 450
%% duplicateId: 0
%% body:
如图，空间某区域中有一匀强磁场，磁感应强度方向水平，且垂直于纸面向里，磁场上边界 b 和下边界 d 水平。在竖直面内有一矩形金属线圈，线圈上下边的距离很短，下边水平。线圈从水平面 a 开始下落。已知磁场上下边界之间的距离大于水平面 a、b 之间的距离。若线圈下边刚通过水平面 b、c（位于磁场中）和 d 时，线圈所受到的磁场力的大小分别为 $F_{b}$、$F_{c}$ 和 $F_{d}$，则\xzanswer{D}
\onepicture[w=4.5cm]{figs/18.png}
\fourchoices[answer=D]
{$F_{d}>F_{c}>F_{b}$}
{$F_{c}<F_{d}<F_{b}$}
{$F_{c}>F_{b}>F_{d}$}
{$F_{c}<F_{b}<F_{d}$}
%% answer: D
%% memo:
\memoanswer{
线圈从 a 运动到 b 做自由落体运动，在 b 点开始进入磁场受到安培力作用 $F_{b}$，由于线圈上下边的距离很短，进入磁场的过程时间很短，进入磁场后，由于磁通量不变，无感应电流产生，不受安培力作用，在 c 处 $F_{c}=0$，但线圈在磁场中受重力作用，做加速运动，出磁场的过程在 d 处受到的安培力比 b 处必然大。故选项 D 正确。
}

\item
%% number: 19
%% paperName: 2010年全国理综Ⅱ第 19 题 6 分
%% typeId: 2
%% type: 多选题
%% chapter: 交流电
%% point: 变压器
%% method: 无
%% score: 6
%% degree: 450
%% duplicateId: 0
%% body:
图中为一理想变压器，其原线圈与一电压有效值不变的交流电源相连：P 为滑动头。现令 P 从均匀密绕的副线圈最底端开始，沿副线圈匀速上滑，直至白炽灯 L 两端的电压等于其额定电压为止。用 $I_{1}$ 表示流过原线圈的电流，$I_{2}$ 表示流过灯泡的电流，$U_{2}$ 表示灯泡两端的电压，$N_{2}$ 表示灯泡消耗的电功率（这里的电流、电压均指有效值：电功率指平均值）。下列 4 个图中，能够正确反映相应物理量的变化趋势的是\xzanswer{BC}
\onepicture[align=right, w=5.5cm]{figs/19a.png}
\fourchoices[answer=BC, ispicture=true, h=3cm]
{figs/19b.png}
{figs/19c.png}
{figs/19d.png}
{figs/19e.png}
%% answer: BC
%% memo:
\memoanswer{
均匀密绕的副线圈最底端开始，沿副线圈匀速上滑，说明副线圈的匝数在均匀增大。根据理想变压器原理 $\frac{U_{1}}{U_{2}}=\frac{n_{1}}{n_{2}}$ 得 $U_{2}=\frac{n_{2}}{n_{1}}U_{1}=\frac{U_{1}}{n_{1}}kt$ 均匀增大（$k$ 为单位时间内增加的匝数），选项 C 正确。灯泡两端电压由零开始增大时，灯泡的电阻也增大，所描绘的伏安特性曲线应该为 B，选项 B 正确、A 错误。灯泡的功率先增大的快（电阻小）后增大的慢（电阻大），选项 D 错误。
}

\item
%% number: 20
%% paperName: 2010年全国理综Ⅱ第 20 题 6 分
%% typeId: 1
%% type: 单选题
%% chapter: 光学
%% point: 全反射
%% method: 无
%% score: 6
%% degree: 650
%% duplicateId: 0
%% body:
频率不同的两束单色光 1 和 2 以相同的入射角从同一点射入一厚玻璃板后，其光路如右图所示，下列说法正确的是\xzanswer{AD}
\onepicture[align=right, w=4.5cm]{figs/20.png}
\fourchoices[answer=AD]
{单色光 1 的波长小于单色光 2 的波长}
{在玻璃中单色光 1 的传播速度大于单色光 2 的传播速度}
{单色光 1 通过玻璃板所需的时间小于单色光 2 通过玻璃板所需的时间}
{单色光 1 从玻璃到空气的全反射临界角小于单色光 2 从玻璃到空气的全反射临界角}
%% answer: AD
%% memo:
\memoanswer{
由题图可知，1 的光线折射率大，频率大，波长小。在介质中传播速度小，因而产生全反射的临界角小。选项 AD 正确，B 错。由 $n=\frac{\sin i}{\sin\gamma}$，在玻璃中传播的距离 $l=\frac{d}{\cos\gamma}$，传播的速度 $v=\frac{c}{n}$，所以光在玻璃中传播的时间 $t=\frac{l}{v}=\frac{d\sin i}{c\sin\gamma\cos\gamma}=\frac{2d\sin i}{c\sin2\gamma}$，1 光线的折射角小，所经历的时间应该长，选项 C 错误。
}

\item
%% number: 21
%% paperName: 2010年全国理综Ⅱ第 21 题 6 分
%% typeId: 1
%% type: 单选题
%% chapter: 曲线运动
%% point: 人造地球卫星
%% method: 无
%% score: 6
%% degree: 450
%% duplicateId: 0
%% body:
已知地球同步卫星离地面的高度约为地球半径的 6 倍。若某行星的平均密度为地球平均密度的一半，它的同步卫星距其表面的高度是其半径的 2.5 倍，则该行星的自转周期约为\xzanswer{B}
\fourchoices[answer=B]
{6 小时}
{12 小时}
{24 小时}
{36 小时}
%% answer: B
%% memo:
\memoanswer{
地球的同步卫星的周期为 $T_{1}=24$ 小时，轨道半径为 $r_{1}=7R_{1}$，密度 $\rho_{1}$。某行星的同步卫星周期为 $T_{2}$，轨道半径为 $r_{2}=3.5R_{2}$，密度 $\rho_{2}$。根据牛顿第二定律和万有引力定律分别有

$\frac{Gm_{1}\times\rho_{1}\times\frac{4}{3}\pi R_{1}^{3}}{r_{1}^{2}}=m_{1}\left(\frac{2\pi}{T_{1}}\right)^{2}r_{1}$，$\frac{Gm_{2}\times\rho_{2}\times\frac{4}{3}\pi R_{2}^{3}}{r_{2}^{2}}=m_{2}\left(\frac{2\pi}{T_{2}}\right)^{2}r_{2}$

两式化简得 $T_{2}=T_{1}/2=12$ 小时，选项 B 正确。
}

\item
%% number: 22
%% paperName: 2010年全国理综Ⅱ第 22 题 5 分
%% typeId: 4
%% type: 实验题
%% chapter: 直线运动
%% point: 自由落体运动
%% method: 无
%% score: 5
%% degree: 650
%% duplicateId: 0
%% body:
利用图中所示的装置可以研究自由落体运动。实验中需要调整好仪器，接通打点计时器的电源，松开纸带，使重物下落。打点计时器会在纸带上打出一系列的小点。

\twopicture[align=right, showlabel=false, hA=5cm, hB=5cm]{figs/22a.png}{figs/22b.png}

\begin{enumerate}
	\item
	为了测试重物下落的加速度，还需要的实验器材有\tkanswer[1.5cm]{C}。（填入正确选项前的字母）

	（A）天平　　（B）秒表　　（C）米尺
	\item
	若实验中所得到的重物下落的加速度值小于当地的重物加速度值，而实验操作与数据处理均无错误，写出一个你认为可能引起此错误差的原因：\tkanswer[4.5cm]{打点计时器与纸带间的摩擦}。
\end{enumerate}
%% answer: 
\jdanswer{
\begin{enumerate}
	\item
	C

	\item
	打点计时器与纸带间的摩擦

\end{enumerate}
}
%% memo:
\memoanswer{
（1）时间由打点计时器确定，用米尺测定位移。答案 C。

（2）打点计时器与纸带间的摩擦。
}

\item
%% number: 23
%% paperName: 2010年全国理综Ⅱ第 23 题 13 分
%% typeId: 4
%% type: 实验题
%% chapter: 电路
%% point: 电路实验
%% method: 等效替代
%% score: 13
%% degree: 450
%% duplicateId: 0
%% body:
如图，一热敏电阻 $R_{T}$ 放在控温容器 M 内：A 为毫安表，量程 $6\UmA$，内阻为数十欧姆；$E$ 为直流电源，电动势约为 $3\UV$，内阻很小；$R$ 为电阻箱，最大阻值为 $999.9\UO$；S 为开关。已知 $R_{T}$ 在 $95\UdC$ 时阻值为 $150\UO$，在 $20\UdC$ 时的阻值约为 $550\UO$。现要求在降温过程中测量在 $95\UdC\sim20\UdC$ 之间的多个温度下 $R_{T}$ 的阻值。

\onepicture[align=right, w=6cm]{figs/23.png}

\begin{enumerate}
	\item
	在图中画出连线，完成实验原理电路图。
	\item
	完成下列实验步骤中的填空：

	\begin{steps}
		\item
		依照实验原理电路图连线
		\item
		调节控温容器 M 内的温度，使得 $R_{T}$ 温度为 $95\UdC$
		\item
		将电阻箱调到适当的初值，以保证仪器安全
		\item
		闭合开关。调节电阻箱，记录电流表的示数 $I_{0}$，并记录\tkanswer[3cm]{电阻箱的读数 $R_{0}$}。
		\item
		将 $R_{T}$ 的温度降为 $T_{1}$（$20\UdC<T_{1}<95\UdC$）；调节电阻箱，使得电流表的读数\tkanswer[1.5cm]{仍为 $I_{0}$}，记录\tkanswer[3cm]{电阻箱的读数为 $R_{1}$}。
		\item
		温度为 $T_{1}$ 时热敏电阻的电阻值 $R_{T}=$\tkanswer[3.5cm]{$R_{0}-R_{1}+150\UO$}。
		\item
		逐步降低 $T_{1}$ 的数值，直至 $20\UdC$ 为止；在每一温度下重复步骤⑤⑥。
	\end{steps}
\end{enumerate}
%% answer: 
\jdanswer{
\begin{enumerate}
	\item
	电路图如图所示。

	\item
	④电阻箱的读数 $R_{0}$；⑤仍为 $I_{0}$，电阻箱的读数为 $R_{1}$；⑥$R_{0}-R_{1}+150\UO$。

\end{enumerate}
}
%% memo:
\memoanswer{
（1）由于本实验只有一只可以测量和观察的直流电流表，所以应该用“替代法”，考虑到用电流表观察而保证电路中电阻不变，因此将热敏电阻、电阻箱和电流表串联形成测量电路。而且热敏电阻 $R_{T}$ 阻值在 $95\UdC$ 和 $20\UdC$ 是已知的，所以热敏电阻的初始温度为 $95\UdC$，则电流表示数不变时，电阻箱和热敏电阻的阻值应保持 $150\UO$ 和电阻箱的初值之和不变。如果热敏电阻的初始温度为 $20\UdC$，则电流表示数不变时，电阻箱和热敏电阻的阻值应保持 $550\UO$ 和电阻箱的初值之和不变。则可以测量任意温度下的电阻。电路图如图所示。

\onepicture[w=6cm]{figs/23_an.png}

（2）①依照实验原理电路图连线

②调节控温容器 M 内的温度，使得 $R_{T}$ 温度为 $95\UdC$

③将电阻箱调到适当的初值，以保证仪器安全

④闭合开关。调节电阻箱，记录电流表的示数 $I_{0}$，并记录电阻箱的读数 $R_{0}$。

⑤将 $R_{T}$ 的温度降为 $T_{1}$（$20\UdC<T_{1}<95\UdC$）；调节电阻箱，使得电流表的读数仍为 $I_{0}$，记录电阻箱的读数为 $R_{1}$。

⑥温度为 $T_{1}$ 时热敏电阻的电阻值 $R_{T}=R_{0}-R_{1}+150\UO$。

⑦逐步降低 $T_{1}$ 的数值，直至 $20\UdC$ 为止；在每一温度下重复步骤⑤⑥即可。
}

\item
%% number: 24
%% paperName: 2010年全国理综Ⅱ第 24 题 15 分
%% typeId: 6
%% type: 计算题
%% chapter: 功和能
%% point: 动能定理求路程
%% method: 无
%% score: 15
%% degree: 450
%% duplicateId: 0
%% body:
如图，MNP 为竖直面内一固定轨道，其圆弧段 MN 与水平段 NP 相切于 N、P 端固定一竖直挡板。M 相对于 N 的高度为 $h$，NP 长度为 $s$。一木块自 M 端从静止开始沿轨道下滑，与挡板发生一次完全弹性碰撞后停止在水平轨道上某处。若在 MN 段的摩擦可忽略不计，物块与 NP 段轨道间的滑动摩擦因数为 $\mu$，求物块停止的地方与 N 点距离的可能值。
\onepicture[align=right, w=6cm]{figs/24.png}
%% answer: 
\jdanswer{
$2s-\frac{h}{\mu}$ 或 $\frac{h}{\mu}-2s$
}
%% memo:
\memoanswer{
根据功能原理，在物块从开始下滑到停止在水平轨道上的过程中，物块的重力势能的减少 $\Delta E_{p}$ 与物块克服摩擦力所做功的数值相等。

$\Delta E_{p}=W$ ①

设物块的质量为 $m$，在水平轨道上滑行的总路程为 $s'$，则

$\Delta E_{p}=mgh$ ②

$W=\mu mgs'$ ③

联立①②③化简得

$s'=\frac{h}{\mu}$ ④

第一种可能是：物块与弹性挡板碰撞后，在 N 前停止，则物块停止的位置距 N 的距离为

$d=2s-s'=2s-\frac{h}{\mu}$ ⑤

第二种可能是：物块与弹性挡板碰撞后，可再一次滑上光滑圆弧轨道，滑下后在水平轨道上停止，则物块停止的位置距 N 的距离为

$d=s'-2s=\frac{h}{\mu}-2s$ ⑥

所以物块停止的位置距 N 的距离可能为 $2s-\frac{h}{\mu}$ 或 $\frac{h}{\mu}-2s$。
}

\item
%% number: 25
%% paperName: 2010年全国理综Ⅱ第 25 题 18 分
%% typeId: 1
%% type: 单选题
%% chapter: 运动和力
%% point: 动量和能量
%% method: 无
%% score: 18
%% degree: 850
%% duplicateId: 0
%% body:
小球 A 和 B 的质量分别为 $m_{A}$ 和 $m_{B}$ 且 $m_{A}>m_{B}$，在某高度处将 A 和 B 先后从静止释放。小球 A 与水平地面碰撞后向上弹回，在释放处的下方与释放初距离为 $H$ 的地方恰好与正在下落的小球 B 发生正碰，设所有碰撞都是弹性的，碰撞时间极短。求小球 A、B 碰撞后 B 上升的最大高度。
%% answer: ABABH
%% memo:
\memoanswer{
根据题意，由运动学规律可知，小球 A 与 B 碰撞前的速度大小相等，设均为 $v_{0}$，由机械能守恒有

$m_{A}gH=\frac{1}{2}m_{A}v_{0}^{2}$ ①

设小球 A 与 B 碰撞后的速度分别为 $v_{1}$ 和 $v_{2}$，以竖直向上方向为正，由动量守恒有

$m_{A}v_{0}+m_{B}(-v_{0})=m_{A}v_{1}+m_{B}v_{2}$ ②

由于两球碰撞过程中能量守恒，故

$\frac{1}{2}m_{A}v_{0}^{2}+\frac{1}{2}m_{B}v_{0}^{2}=\frac{1}{2}m_{A}v_{1}^{2}+\frac{1}{2}m_{B}v_{2}^{2}$ ③

联立②③式得

$v_{2}=\frac{3m_{A}-m_{B}}{m_{A}+m_{B}}v_{0}$ ④

设小球 B 能上升的最大高度为 $h$，由运动学公式有

$h=\frac{v_{2}^{2}}{2g}$ ⑤

由①④⑤式得

$h=\left(\frac{3m_{A}-m_{B}}{m_{A}+m_{B}}\right)^{2}H$ ⑥
}

\item
%% number: 26
%% paperName: 2010年全国理综Ⅱ第 26 题 21 分
%% typeId: 6
%% type: 计算题
%% chapter: 磁场
%% point: 带电粒子在磁场中的运动
%% method: 无
%% score: 21
%% degree: 250
%% duplicateId: 0
%% body:
图中左边有一对平行金属板，两板相距为 $d$，电压为 $U$；两板之间有匀强磁场，磁场应强度大小为 $B_{0}$，方向平行于板面并垂直于纸面朝里。图中右边有一边长为 $a$ 的正三角形区域 EFG（EF 边与金属板垂直），在此区域内及其边界上也有匀强磁场，磁感应强度大小为 $B$，方向垂直于纸面朝里。假设一系列电荷量为 $q$ 的正离子沿平行于金属板面，垂直于磁场的方向射入金属板之间，沿同一方向射出金属板之间的区域，并经 EF 边中点 H 射入磁场区域。不计重力。

\begin{enumerate}
	\item
	已知这些离子中的离子甲到达磁场边界 EG 后，从边界 EF 穿出磁场，求离子甲的质量。
	\item
	已知这些离子中的离子乙从 EG 边上的 I 点（图中未画出）穿出磁场，且 GI 长为 $\frac{3}{4}a$，求离子乙的质量。
	\item
	若这些离子中的最轻离子的质量等于离子甲质量的一半，而离子乙的质量是最大的，问磁场边界上什么区域内可能有离子到达。
\end{enumerate}
\onepicture[align=right, w=7cm]{figs/26.png}
%% answer: 
\jdanswer{
\begin{enumerate}
	\item
	$m=\frac{qaBB_{0}d}{U}\left(\sqrt{3}-\frac{3}{2}\right)$

	\item
	$m'=\frac{qaBB_{0}d}{4U}$

	\item
	EF 边上从 O 到 P，EG 边上从 K 到 I

\end{enumerate}
}
%% memo:
\memoanswer{
（1）由题意知，所有离子在平行金属板之间做匀速直线运动，它所受到的向上的磁场力和向下的电场力平衡，有

$qvB_{0}=qE_{0}$ ①

式中 $v$ 是离子运动的速度，$E_{0}$ 是平行金属板之间的匀强电场的强度，有

$E_{0}=\frac{U}{d}$ ②

由①②式得

$v=\frac{U}{B_{0}d}$ ③

在正三角形磁场区域，离子甲做匀速圆周运动。设离子甲质量为 $m$，由洛伦兹力公式和牛顿第二定律有

$qvB=m\frac{v^{2}}{r}$ ④

式中，$r$ 是离子甲做圆周运动的半径。离子甲在磁场中的运动轨迹为半圆，圆心为 O：这半圆刚好与 EG 边相切于 K，与 EF 边交于 I 点。在 $\triangle EOK$ 中，OK 垂直于 EG。由几何关系得

$\frac{1}{2}a-r=\frac{2}{\sqrt{3}}r$ ⑤

\onepicture[w=6cm]{figs/26_an1.png}

由⑤式得

$r=\left(\sqrt{3}-\frac{3}{2}\right)a$ ⑥

联立③④⑥式得，离子甲的质量为

$m=\frac{qaBB_{0}d}{U}\left(\sqrt{3}-\frac{3}{2}\right)$ ⑦

（2）同理，有洛伦兹力公式和牛顿第二定律有

$qvB=m'\frac{v^{2}}{r'}$ ⑧

式中，$m'$ 和 $r'$ 分别为离子乙的质量和做圆周运动的轨道半径。离子乙运动的圆周的圆心 $O'$ 必在 E、H 两点之间，又几何关系有

$r'^{2}=\left(a-\frac{3}{4}a\right)^{2}+\left(\frac{a}{2}-r'\right)^{2}-2\left(a-\frac{3}{4}a\right)\left(\frac{a}{2}-r'\right)\cos60^\circ$ ⑨

\onepicture[w=10cm]{figs/26_an2.png}

由⑨式得

$r'=\frac{1}{4}a$ \ccnum{⑩}

联立③⑧⑩式得，离子乙的质量为

$m'=\frac{qaBB_{0}d}{4U}$ \ccnum{(11)}

（3）对于最轻的离子，其质量为 $m/2$，由④式知，它在磁场中做半径为 $r/2$ 的匀速圆周运动。因而与 EH 的交点为 O，有

$OH=\left(\sqrt{3}-\frac{3}{2}\right)a$ \ccnum{(12)}

\onepicture[w=5cm]{figs/26_an3.png}

当这些离子中的离子质量逐渐增大到 $m$ 时，离子到达磁场边界上的点的位置从 O 点沿 HE 边变到 P 点；当离子质量继续增大时，离子到达磁场边界上的点的位置从 K 点沿 EG 边趋向于 I 点。K 点到 G 点的距离为

$KG=\frac{\sqrt{3}}{2}a$ \ccnum{(13)}

所以，磁场边界上可能有离子到达的区域是：EF 边上从 O 到 P，EG 边上从 K 到 I。
}

\end{enumerate}

\end{document}
